\documentclass[11pt,a4paper]{article}

% ============================================================
% PACKAGES
% ============================================================
\usepackage[margin=1in]{geometry}
\usepackage{hyperref}
\usepackage{listings}
\usepackage{xcolor}
\usepackage{graphicx}
\usepackage{enumitem}
\usepackage{amsmath}
\usepackage{booktabs}
\usepackage{fancyhdr}
\usepackage{tcolorbox}

% ============================================================
% HYPERREF SETUP
% ============================================================
\hypersetup{
    colorlinks=true,
    linkcolor=blue,
    filecolor=magenta,
    urlcolor=cyan,
    pdftitle={Pipeline Parallelism Neural Network with MPI},
    pdfauthor={Technical Documentation},
}

% ============================================================
% LISTINGS CONFIGURATION
% ============================================================
\definecolor{codegreen}{rgb}{0,0.6,0}
\definecolor{codegray}{rgb}{0.5,0.5,0.5}
\definecolor{codepurple}{rgb}{0.58,0,0.82}
\definecolor{backcolour}{rgb}{0.95,0.95,0.92}
\definecolor{keywordcolor}{rgb}{0.0,0.0,0.7}

\lstdefinestyle{mystyle}{
    backgroundcolor=\color{backcolour},
    commentstyle=\color{codegreen},
    keywordstyle=\color{keywordcolor}\bfseries,
    numberstyle=\tiny\color{codegray},
    stringstyle=\color{codepurple},
    basicstyle=\ttfamily\footnotesize,
    breakatwhitespace=false,
    breaklines=true,
    captionpos=b,
    keepspaces=true,
    numbers=left,
    numbersep=5pt,
    showspaces=false,
    showstringspaces=false,
    showtabs=false,
    tabsize=4,
    frame=single,
    framerule=0.5pt,
    rulecolor=\color{codegray},
}

\lstset{style=mystyle}

% ============================================================
% DOCUMENT INFO
% ============================================================
\title{%
    \textbf{Pipeline Parallelism for Fully Connected Neural Networks} \\
    \large Using MPI (Message Passing Interface) \\
    \vspace{0.5cm}
    \normalsize Technical Documentation for \texttt{fcn\_mpi.c}
}
\author{High Performance Computing Project}
\date{February 2026}

% ============================================================
% DOCUMENT BEGIN
% ============================================================
\begin{document}

\maketitle
\tableofcontents
\newpage

% ============================================================
% EXECUTIVE SUMMARY
% ============================================================
\section{Executive Summary}

This document explains a C program that implements \textbf{pipeline parallelism} for training a fully connected neural network (FCN) on the MNIST handwritten digit dataset using \textbf{MPI} (Message Passing Interface).

\begin{tcolorbox}[colback=blue!5!white,colframe=blue!75!black,title=Key Points]
\begin{itemize}[noitemsep]
    \item \textbf{Purpose}: Train a 3-layer neural network using distributed computing across 3 MPI processes.
    \item \textbf{Architecture}: 784 $\rightarrow$ 256 $\rightarrow$ 128 $\rightarrow$ 10 neurons.
    \item \textbf{Parallelism Model}: Pipeline parallelism where each process handles specific layers.
    \item \textbf{Dataset}: MNIST (60,000 training images, 10,000 test images).
    \item \textbf{Communication}: MPI point-to-point operations (\texttt{MPI\_Send}/\texttt{MPI\_Recv}).
    \item \textbf{Training}: Stochastic Gradient Descent (SGD) with backpropagation.
\end{itemize}
\end{tcolorbox}

The program distributes the neural network layers across three processes:
\begin{itemize}
    \item \textbf{Process 0}: Handles input data and Hidden Layer 1
    \item \textbf{Process 1}: Handles Hidden Layer 2
    \item \textbf{Process 2}: Handles Output Layer and loss calculation
\end{itemize}

% ============================================================
% CONTEXT & ASSUMPTIONS
% ============================================================
\section{Context \& Assumptions}

\subsection{Runtime Environment}

\begin{itemize}
    \item \textbf{Operating System}: Linux/Unix-based systems (tested on Linux)
    \item \textbf{Compiler}: \texttt{mpicc} (MPI C compiler wrapper)
    \item \textbf{MPI Implementation}: OpenMPI, MPICH, or compatible MPI library
    \item \textbf{C Standard}: C99 or later
    \item \textbf{Required Libraries}: 
    \begin{itemize}
        \item \texttt{math.h} (mathematical functions)
        \item \texttt{mpi.h} (MPI functions)
        \item Standard C libraries (\texttt{stdio.h}, \texttt{stdlib.h}, \texttt{string.h}, \texttt{stdint.h}, \texttt{time.h})
    \end{itemize}
\end{itemize}

\subsection{Input Requirements}

The program expects MNIST dataset files in the \texttt{data/} directory:

\begin{table}[h]
\centering
\begin{tabular}{@{}ll@{}}
\toprule
\textbf{File} & \textbf{Description} \\
\midrule
\texttt{train-images-idx3-ubyte} & Training images (60,000 samples) \\
\texttt{train-labels-idx1-ubyte} & Training labels \\
\texttt{t10k-images-idx3-ubyte} & Test images (10,000 samples) \\
\texttt{t10k-labels-idx1-ubyte} & Test labels \\
\bottomrule
\end{tabular}
\caption{Required MNIST dataset files}
\end{table}

\subsection{Output}

\begin{itemize}
    \item Training progress: Loss per epoch and time per epoch
    \item Final test accuracy percentage
    \item Total execution time
\end{itemize}

\subsection{Assumptions}

\begin{enumerate}
    \item Exactly 3 MPI processes are launched
    \item MNIST files are in IDX format (standard MNIST binary format)
    \item Sufficient memory available on each process for layer weights and activations
    \item All processes can communicate via MPI (shared memory or network)
\end{enumerate}

% ============================================================
% HIGH-LEVEL ARCHITECTURE
% ============================================================
\section{High-Level Architecture}

\subsection{Module Organization}

The code is organized into logical sections within a single file:

\begin{figure}[h]
\centering
\begin{tcolorbox}[colback=gray!10,colframe=gray!50,width=0.9\textwidth]
\begin{verbatim}
fcn_mpi.c
├── Configuration (Macros/Constants)
├── Data Structures
│   ├── Tensor
│   └── DenseLayer
├── Tensor Operations
├── Dense Layer Operations
├── Activation Functions (ReLU, Softmax)
├── Loss Functions (Cross-Entropy)
├── Dataset Loading (MNIST IDX format)
├── Process-specific Functions
│   ├── process_0_run() - Input + Layer 1
│   ├── process_1_run() - Layer 2
│   └── process_2_run() - Output + Loss
└── Main Function
\end{verbatim}
\end{tcolorbox}
\caption{Code organization}
\end{figure}

\subsection{Pipeline Parallelism Model}

\textbf{Pipeline parallelism} divides a neural network vertically by layers, with each process responsible for one or more consecutive layers.

\begin{figure}[h]
\centering
\begin{tcolorbox}[colback=green!5,colframe=green!50!black,width=0.95\textwidth]
\textbf{Forward Pass Data Flow:}
\begin{verbatim}
Process 0              Process 1              Process 2
─────────              ─────────              ─────────
Input (784)
    │
Dense Layer 1
(784 → 256)
    │
ReLU
    │
    └──── a1 (256) ────►
                       Dense Layer 2
                       (256 → 128)
                           │
                         ReLU
                           │
                           └──── a2 (128) ────►
                                               Dense Layer 3
                                               (128 → 10)
                                                   │
                                               Softmax
                                                   │
                                               Loss ◄── label
\end{verbatim}
\end{tcolorbox}
\caption{Forward pass data flow between processes}
\end{figure}

\begin{figure}[h]
\centering
\begin{tcolorbox}[colback=red!5,colframe=red!50!black,width=0.95\textwidth]
\textbf{Backward Pass Gradient Flow:}
\begin{verbatim}
Process 0              Process 1              Process 2
─────────              ─────────              ─────────
                                               Loss Gradient
                                                   │
                                               Dense3 Backward
                                                   │
                       ◄──── grad_a2 (128) ────────┘
                       ReLU Backward
                           │
                       Dense2 Backward
                           │
◄──── grad_a1 (256) ───────┘
ReLU Backward
    │
Dense1 Backward
    │
(No further propagation)
\end{verbatim}
\end{tcolorbox}
\caption{Backward pass gradient flow between processes}
\end{figure}

\subsection{Process Responsibilities}

\begin{table}[h]
\centering
\begin{tabular}{@{}clll@{}}
\toprule
\textbf{Process} & \textbf{Layers} & \textbf{Dimensions} & \textbf{Additional Tasks} \\
\midrule
0 & Input + Dense1 + ReLU & 784 $\rightarrow$ 256 & Load dataset, print results \\
1 & Dense2 + ReLU & 256 $\rightarrow$ 128 & Relay gradients \\
2 & Dense3 + Softmax & 128 $\rightarrow$ 10 & Calculate loss, accuracy \\
\bottomrule
\end{tabular}
\caption{Process responsibilities in pipeline}
\end{table}

% ============================================================
% STEP-BY-STEP WALKTHROUGH
% ============================================================
\section{Step-by-Step Walkthrough}

This section describes the execution flow from program start to finish.

\subsection{Step 1: MPI Initialization}

\begin{lstlisting}[language=C, caption={MPI initialization in main()}]
int main(int argc, char *argv[]) {
    int rank, size;

    MPI_Init(&argc, &argv);
    MPI_Comm_rank(MPI_COMM_WORLD, &rank);
    MPI_Comm_size(MPI_COMM_WORLD, &size);
\end{lstlisting}

\begin{itemize}
    \item \texttt{MPI\_Init}: Initializes the MPI execution environment
    \item \texttt{MPI\_Comm\_rank}: Gets the unique identifier (rank) of the calling process
    \item \texttt{MPI\_Comm\_size}: Gets the total number of processes
\end{itemize}

\subsection{Step 2: Validation and Setup}

\begin{lstlisting}[language=C, caption={Process count validation}]
    if (size != NUM_PROCESSES) {
        if (rank == 0) {
            fprintf(stderr, "Error: This program requires exactly %d MPI processes.\n", NUM_PROCESSES);
        }
        MPI_Finalize();
        return 1;
    }

    srand(42 + rank);  /* Different seed for each process, but reproducible */
\end{lstlisting}

\begin{itemize}
    \item Ensures exactly 3 processes are running
    \item Each process gets a unique but reproducible random seed
\end{itemize}

\subsection{Step 3: Dataset Loading (Process 0 Only)}

\begin{lstlisting}[language=C, caption={Dataset loading by Process 0}]
    if (rank == 0) {
        load_mnist(&train, &test);
        num_train_samples = train.num_samples;
        num_test_samples = test.num_samples;
    }

    /* Broadcast dataset sizes to all processes */
    MPI_Bcast(&num_train_samples, 1, MPI_INT, 0, MPI_COMM_WORLD);
    MPI_Bcast(&num_test_samples, 1, MPI_INT, 0, MPI_COMM_WORLD);
\end{lstlisting}

\begin{itemize}
    \item Only Process 0 loads the MNIST dataset from disk
    \item Dataset sizes are broadcast to all processes using \texttt{MPI\_Bcast}
    \item This is a \textbf{collective operation} -- all processes must call it
\end{itemize}

\subsection{Step 4: Process Dispatch}

\begin{lstlisting}[language=C, caption={Process dispatch based on rank}]
    switch (rank) {
        case 0:
            process_0_run(&train, &test, num_train_samples);
            break;
        case 1:
            process_1_run(num_train_samples, num_test_samples);
            break;
        case 2:
            process_2_run(num_train_samples, num_test_samples);
            break;
    }
\end{lstlisting}

Each process executes its designated function based on its rank.

\subsection{Step 5: Training Loop (Per Sample)}

For each training sample, the following sequence occurs:

\begin{enumerate}
    \item \textbf{Process 0}:
    \begin{enumerate}[label=\alph*.]
        \item Load image from dataset
        \item Forward pass through Dense Layer 1 + ReLU
        \item Send activation \texttt{a1} to Process 1
        \item Send label to Process 2
        \item Wait for gradient from Process 1
        \item Backward pass and weight update
    \end{enumerate}
    
    \item \textbf{Process 1}:
    \begin{enumerate}[label=\alph*.]
        \item Receive activation \texttt{a1} from Process 0
        \item Forward pass through Dense Layer 2 + ReLU
        \item Send activation \texttt{a2} to Process 2
        \item Wait for gradient from Process 2
        \item Backward pass
        \item Send gradient to Process 0
        \item Weight update
    \end{enumerate}
    
    \item \textbf{Process 2}:
    \begin{enumerate}[label=\alph*.]
        \item Receive activation \texttt{a2} from Process 1
        \item Receive label from Process 0
        \item Forward pass through Dense Layer 3 + Softmax
        \item Calculate cross-entropy loss
        \item Backward pass
        \item Send gradient to Process 1
        \item Weight update
    \end{enumerate}
\end{enumerate}

\subsection{Step 6: Epoch Completion}

After all samples in an epoch:
\begin{itemize}
    \item Process 2 calculates average loss and sends to Process 0
    \item Process 0 prints epoch statistics (loss, time)
\end{itemize}

\subsection{Step 7: Evaluation}

After training completes:
\begin{itemize}
    \item Forward pass only (no backpropagation)
    \item Process 2 counts correct predictions
    \item Accuracy sent to Process 0 for display
\end{itemize}

\subsection{Step 8: Cleanup}

\begin{lstlisting}[language=C, caption={Cleanup and finalization}]
    if (rank == 0) {
        printf("\nTotal execution time: %.2f seconds\n", end_time - start_time);
        free_dataset(&train);
        free_dataset(&test);
    }

    MPI_Finalize();
    return 0;
\end{lstlisting}

% ============================================================
% DETAILED EXPLANATION BY COMPONENT
% ============================================================
\section{Detailed Explanation by Component}

\subsection{Configuration Constants}

\begin{lstlisting}[language=C, caption={Configuration macros}]
#define INPUT_SIZE     784      // 28x28 pixel images flattened
#define HIDDEN1_SIZE   256      // First hidden layer neurons
#define HIDDEN2_SIZE   128      // Second hidden layer neurons
#define NUM_CLASSES    10       // Digits 0-9

#define EPOCHS         10       // Training iterations over dataset
#define LEARNING_RATE  0.01f    // SGD step size
#define EPSILON        1e-7f    // Small value to prevent log(0)

#define NUM_PROCESSES  3        // Required MPI processes
\end{lstlisting}

\begin{itemize}
    \item \texttt{INPUT\_SIZE}: MNIST images are 28$\times$28 = 784 pixels
    \item \texttt{HIDDEN1\_SIZE}, \texttt{HIDDEN2\_SIZE}: Hidden layer dimensions (can be tuned)
    \item \texttt{NUM\_CLASSES}: 10 digit classes (0-9)
    \item \texttt{LEARNING\_RATE}: Controls how much weights change per update
    \item \texttt{EPSILON}: Prevents numerical instability in logarithm
\end{itemize}

\subsection{MPI Tags}

\begin{lstlisting}[language=C, caption={MPI communication tags}]
#define TAG_FORWARD_A1     100  // Activation from layer 1
#define TAG_FORWARD_A2     101  // Activation from layer 2
#define TAG_FORWARD_LABEL  102  // Label for loss calculation
#define TAG_BACKWARD_GRAD1 200  // Gradient to layer 1
#define TAG_BACKWARD_GRAD2 201  // Gradient to layer 2
#define TAG_LOSS           300  // Loss/accuracy value
\end{lstlisting}

\textbf{Tags} in MPI distinguish different message types, allowing a process to receive specific messages.

\subsection{Tensor Structure}

\begin{lstlisting}[language=C, caption={Tensor data structure}]
typedef struct {
    float *data;    // Pointer to contiguous float array
    int rows;       // Number of rows
    int cols;       // Number of columns
} Tensor;
\end{lstlisting}

\textbf{Purpose}: A simple 2D matrix representation for storing weights, activations, and gradients.

\subsubsection{tensor\_create}

\begin{lstlisting}[language=C, caption={Tensor creation function}]
Tensor tensor_create(int rows, int cols) {
    Tensor t;
    t.rows = rows;
    t.cols = cols;
    t.data = (float *)calloc(rows * cols, sizeof(float));
    return t;
}
\end{lstlisting}

\begin{itemize}
    \item \textbf{Parameters}: \texttt{rows}, \texttt{cols} -- dimensions of the tensor
    \item \textbf{Returns}: Initialized \texttt{Tensor} struct
    \item \textbf{Memory}: Uses \texttt{calloc} to zero-initialize memory
    \item \textbf{Side Effects}: Allocates heap memory (must be freed)
\end{itemize}

\subsubsection{tensor\_free}

\begin{lstlisting}[language=C, caption={Tensor deallocation}]
void tensor_free(Tensor *t) {
    if (t->data) {
        free(t->data);
        t->data = NULL;
    }
}
\end{lstlisting}

\begin{itemize}
    \item \textbf{Purpose}: Safely deallocates tensor memory
    \item \textbf{Side Effects}: Sets \texttt{data} to \texttt{NULL} to prevent double-free
\end{itemize}

\subsubsection{tensor\_zero}

\begin{lstlisting}[language=C, caption={Zero-fill tensor}]
void tensor_zero(Tensor *t) {
    memset(t->data, 0, t->rows * t->cols * sizeof(float));
}
\end{lstlisting}

\begin{itemize}
    \item \textbf{Purpose}: Sets all elements to zero
    \item \textbf{Usage}: Initializing bias vectors, clearing gradients
\end{itemize}

\subsection{Dense Layer Structure}

\begin{lstlisting}[language=C, caption={Dense layer data structure}]
typedef struct {
    Tensor weights;      // Weight matrix (out_size x in_size)
    Tensor bias;         // Bias vector (out_size x 1)
    Tensor grad_w;       // Weight gradients
    Tensor grad_b;       // Bias gradients
    Tensor input_cache;  // Cached input for backprop
    int in_size;         // Input dimension
    int out_size;        // Output dimension
} DenseLayer;
\end{lstlisting}

A \textbf{dense layer} (fully connected layer) computes: $\mathbf{y} = \mathbf{W}\mathbf{x} + \mathbf{b}$

\subsubsection{dense\_init}

\begin{lstlisting}[language=C, caption={Dense layer initialization}]
void dense_init(DenseLayer *l, int in, int out) {
    l->in_size = in;
    l->out_size = out;
    l->weights = tensor_create(out, in);
    l->bias = tensor_create(out, 1);
    l->grad_w = tensor_create(out, in);
    l->grad_b = tensor_create(out, 1);
    l->input_cache = tensor_create(in, 1);

    for (int i = 0; i < out * in; i++) {
        l->weights.data[i] = rand_uniform(-0.05f, 0.05f);
    }
    tensor_zero(&l->bias);
}
\end{lstlisting}

\begin{itemize}
    \item \textbf{Parameters}: \texttt{in} (input size), \texttt{out} (output size)
    \item \textbf{Weight Initialization}: Random uniform in $[-0.05, 0.05]$
    \item \textbf{Bias Initialization}: Zero
    \item \textbf{Memory}: Allocates 5 tensors per layer
\end{itemize}

\subsubsection{dense\_forward}

\begin{lstlisting}[language=C, caption={Dense layer forward pass}]
void dense_forward(DenseLayer *l, float *input, float *output) {
    int out = l->out_size;
    int in = l->in_size;

    /* Cache input for backward pass */
    memcpy(l->input_cache.data, input, in * sizeof(float));

    for (int i = 0; i < out; i++) {
        float sum = l->bias.data[i];
        for (int j = 0; j < in; j++) {
            sum += l->weights.data[i * in + j] * input[j];
        }
        output[i] = sum;
    }
}
\end{lstlisting}

\textbf{Mathematical Operation}:
\begin{equation}
y_i = b_i + \sum_{j=0}^{\text{in}-1} W_{ij} \cdot x_j
\end{equation}

\begin{itemize}
    \item \textbf{Input Caching}: Stores input for use in backward pass
    \item \textbf{Time Complexity}: $O(\text{in} \times \text{out})$
\end{itemize}

\subsubsection{dense\_backward}

\begin{lstlisting}[language=C, caption={Dense layer backward pass}]
void dense_backward(DenseLayer *l, float *grad_out, float *grad_in) {
    int out = l->out_size;
    int in = l->in_size;

    /* dW = grad_out * input^T */
    for (int i = 0; i < out; i++) {
        for (int j = 0; j < in; j++) {
            l->grad_w.data[i * in + j] = grad_out[i] * l->input_cache.data[j];
        }
    }

    /* db = grad_out */
    for (int i = 0; i < out; i++) {
        l->grad_b.data[i] = grad_out[i];
    }

    /* dInput = W^T * grad_out */
    if (grad_in != NULL) {
        for (int j = 0; j < in; j++) {
            float sum = 0.0f;
            for (int i = 0; i < out; i++) {
                sum += l->weights.data[i * in + j] * grad_out[i];
            }
            grad_in[j] = sum;
        }
    }
}
\end{lstlisting}

\textbf{Gradient Calculations} (using chain rule):
\begin{align}
\frac{\partial L}{\partial W_{ij}} &= \frac{\partial L}{\partial y_i} \cdot x_j \\
\frac{\partial L}{\partial b_i} &= \frac{\partial L}{\partial y_i} \\
\frac{\partial L}{\partial x_j} &= \sum_{i=0}^{\text{out}-1} W_{ij} \cdot \frac{\partial L}{\partial y_i}
\end{align}

\begin{itemize}
    \item \texttt{grad\_out}: Gradient from the next layer ($\frac{\partial L}{\partial y}$)
    \item \texttt{grad\_in}: Gradient to propagate to previous layer (can be \texttt{NULL} for first layer)
\end{itemize}

\subsubsection{sgd\_update}

\begin{lstlisting}[language=C, caption={SGD weight update}]
void sgd_update(DenseLayer *l, float lr) {
    int n_w = l->out_size * l->in_size;
    for (int i = 0; i < n_w; i++) {
        l->weights.data[i] -= lr * l->grad_w.data[i];
    }
    for (int i = 0; i < l->out_size; i++) {
        l->bias.data[i] -= lr * l->grad_b.data[i];
    }
}
\end{lstlisting}

\textbf{SGD Update Rule}:
\begin{equation}
\theta_{t+1} = \theta_t - \eta \cdot \nabla_\theta L
\end{equation}

where $\eta$ is the learning rate and $\nabla_\theta L$ is the gradient.

\subsection{Activation Functions}

\subsubsection{ReLU (Rectified Linear Unit)}

\begin{lstlisting}[language=C, caption={ReLU activation}]
void relu_forward(float *x, float *out, int n) {
    for (int i = 0; i < n; i++) {
        out[i] = (x[i] > 0.0f) ? x[i] : 0.0f;
    }
}

void relu_backward(float *x, float *grad_out, float *grad_in, int n) {
    for (int i = 0; i < n; i++) {
        grad_in[i] = (x[i] > 0.0f) ? grad_out[i] : 0.0f;
    }
}
\end{lstlisting}

\textbf{ReLU Function}:
\begin{equation}
\text{ReLU}(x) = \max(0, x)
\end{equation}

\textbf{ReLU Derivative}:
\begin{equation}
\frac{\partial \text{ReLU}}{\partial x} = 
\begin{cases}
1 & \text{if } x > 0 \\
0 & \text{otherwise}
\end{cases}
\end{equation}

\subsubsection{Softmax}

\begin{lstlisting}[language=C, caption={Softmax activation}]
void softmax(float *x, float *out, int n) {
    float max_val = x[0];
    for (int i = 1; i < n; i++) {
        if (x[i] > max_val) max_val = x[i];
    }

    float sum = 0.0f;
    for (int i = 0; i < n; i++) {
        out[i] = expf(x[i] - max_val);  // Subtract max for numerical stability
        sum += out[i];
    }

    for (int i = 0; i < n; i++) {
        out[i] /= sum;
    }
}
\end{lstlisting}

\textbf{Softmax Function}:
\begin{equation}
\text{softmax}(x_i) = \frac{e^{x_i - \max(x)}}{\sum_{j} e^{x_j - \max(x)}}
\end{equation}

The $\max(x)$ subtraction prevents numerical overflow when exponentiating large values.

\subsection{Loss Functions}

\subsubsection{Cross-Entropy Loss}

\begin{lstlisting}[language=C, caption={Cross-entropy loss}]
float cross_entropy_loss(float *probs, int label) {
    float p = probs[label];
    if (p < EPSILON) p = EPSILON;  // Prevent log(0)
    return -logf(p);
}
\end{lstlisting}

\textbf{Cross-Entropy Loss}:
\begin{equation}
L = -\log(p_{\text{label}})
\end{equation}

where $p_{\text{label}}$ is the predicted probability for the correct class.

\subsubsection{Cross-Entropy Backward}

\begin{lstlisting}[language=C, caption={Cross-entropy gradient}]
void cross_entropy_backward(float *probs, int label, float *grad, int n) {
    for (int i = 0; i < n; i++) {
        grad[i] = probs[i];
    }
    grad[label] -= 1.0f;
}
\end{lstlisting}

\textbf{Gradient} (combined softmax + cross-entropy):
\begin{equation}
\frac{\partial L}{\partial z_i} = p_i - \mathbf{1}_{i=\text{label}}
\end{equation}

This elegant formula is the result of combining the softmax and cross-entropy derivatives.

\subsection{Dataset Loading}

\subsubsection{IDX File Format}

MNIST uses the IDX file format:

\begin{tcolorbox}[colback=yellow!10,colframe=yellow!50!black,title=IDX Image File Format]
\begin{verbatim}
[offset] [type]          [description]
0000     32-bit integer  magic number (0x00000803)
0004     32-bit integer  number of images
0008     32-bit integer  number of rows (28)
0012     32-bit integer  number of columns (28)
0016     unsigned byte   pixel data (row-major)
...
\end{verbatim}
\end{tcolorbox}

\subsubsection{read\_u32}

\begin{lstlisting}[language=C, caption={Big-endian 32-bit read}]
static uint32_t read_u32(FILE *f) {
    uint8_t b[4];
    fread(b, 1, 4, f);
    return (b[0] << 24) | (b[1] << 16) | (b[2] << 8) | b[3];
}
\end{lstlisting}

IDX files use \textbf{big-endian} byte order. This function reads 4 bytes and converts to native integer.

\subsubsection{read\_idx\_images}

\begin{lstlisting}[language=C, caption={Reading MNIST images}]
float *read_idx_images(const char *path, int *count) {
    FILE *f = fopen(path, "rb");
    if (!f) {
        fprintf(stderr, "Error: Cannot open %s\n", path);
        return NULL;
    }

    read_u32(f);               /* magic number - ignored */
    *count = read_u32(f);      /* number of images */
    int rows = read_u32(f);    /* 28 */
    int cols = read_u32(f);    /* 28 */

    int size = (*count) * rows * cols;
    float *data = (float *)malloc(sizeof(float) * size);

    for (int i = 0; i < size; i++) {
        unsigned char pixel;
        fread(&pixel, 1, 1, f);
        data[i] = pixel / 255.0f;  // Normalize to [0, 1]
    }

    fclose(f);
    return data;
}
\end{lstlisting}

\begin{itemize}
    \item Reads binary pixel data
    \item Normalizes from $[0, 255]$ to $[0.0, 1.0]$
    \item Returns flattened array (all images concatenated)
\end{itemize}

\subsection{Process Functions}

\subsubsection{process\_0\_run (Input + Layer 1)}

\begin{lstlisting}[language=C, caption={Process 0 main loop (training)}]
void process_0_run(Dataset *train, Dataset *test, int num_samples) {
    DenseLayer l1;
    dense_init(&l1, INPUT_SIZE, HIDDEN1_SIZE);

    /* Allocate buffers */
    float *x = (float *)malloc(INPUT_SIZE * sizeof(float));
    float *z1 = (float *)malloc(HIDDEN1_SIZE * sizeof(float));
    float *a1 = (float *)malloc(HIDDEN1_SIZE * sizeof(float));
    float *grad_a1 = (float *)malloc(HIDDEN1_SIZE * sizeof(float));
    float *grad_z1 = (float *)malloc(HIDDEN1_SIZE * sizeof(float));

    for (int epoch = 0; epoch < EPOCHS; epoch++) {
        for (int i = 0; i < num_samples; i++) {
            /* Load input */
            memcpy(x, &train->images[i * INPUT_SIZE], INPUT_SIZE * sizeof(float));
            int label = train->labels[i];

            /* Forward pass */
            dense_forward(&l1, x, z1);
            relu_forward(z1, a1, HIDDEN1_SIZE);

            /* Send to other processes */
            MPI_Send(a1, HIDDEN1_SIZE, MPI_FLOAT, 1, TAG_FORWARD_A1, MPI_COMM_WORLD);
            MPI_Send(&label, 1, MPI_INT, 2, TAG_FORWARD_LABEL, MPI_COMM_WORLD);

            /* Receive gradient */
            MPI_Recv(grad_a1, HIDDEN1_SIZE, MPI_FLOAT, 1, TAG_BACKWARD_GRAD1, 
                     MPI_COMM_WORLD, &status);

            /* Backward pass */
            relu_backward(z1, grad_a1, grad_z1, HIDDEN1_SIZE);
            dense_backward(&l1, grad_z1, NULL);

            /* Update weights */
            sgd_update(&l1, LEARNING_RATE);
        }
        /* ... epoch statistics ... */
    }
}
\end{lstlisting}

\textbf{Key MPI Operations}:
\begin{itemize}
    \item \texttt{MPI\_Send(a1, ...)}: Sends 256 floats to Process 1
    \item \texttt{MPI\_Send(\&label, ...)}: Sends integer label to Process 2
    \item \texttt{MPI\_Recv(grad\_a1, ...)}: Blocking receive for gradient
\end{itemize}

\subsubsection{process\_1\_run (Layer 2)}

Process 1 acts as a \textbf{relay} between Process 0 and Process 2:

\begin{lstlisting}[language=C, caption={Process 1 main loop}]
void process_1_run(int num_train_samples, int num_test_samples) {
    DenseLayer l2;
    dense_init(&l2, HIDDEN1_SIZE, HIDDEN2_SIZE);

    for (int epoch = 0; epoch < EPOCHS; epoch++) {
        for (int i = 0; i < num_train_samples; i++) {
            /* Receive from Process 0 */
            MPI_Recv(a1, HIDDEN1_SIZE, MPI_FLOAT, 0, TAG_FORWARD_A1, ...);

            /* Forward pass */
            dense_forward(&l2, a1, z2);
            relu_forward(z2, a2, HIDDEN2_SIZE);

            /* Send to Process 2 */
            MPI_Send(a2, HIDDEN2_SIZE, MPI_FLOAT, 2, TAG_FORWARD_A2, ...);

            /* Receive gradient from Process 2 */
            MPI_Recv(grad_a2, HIDDEN2_SIZE, MPI_FLOAT, 2, TAG_BACKWARD_GRAD2, ...);

            /* Backward pass */
            relu_backward(z2, grad_a2, grad_z2, HIDDEN2_SIZE);
            dense_backward(&l2, grad_z2, grad_a1);

            /* Send gradient to Process 0 */
            MPI_Send(grad_a1, HIDDEN1_SIZE, MPI_FLOAT, 0, TAG_BACKWARD_GRAD1, ...);

            /* Update weights */
            sgd_update(&l2, LEARNING_RATE);
        }
    }
}
\end{lstlisting}

\subsubsection{process\_2\_run (Output + Loss)}

\begin{lstlisting}[language=C, caption={Process 2 main loop}]
void process_2_run(int num_train_samples, int num_test_samples) {
    DenseLayer l3;
    dense_init(&l3, HIDDEN2_SIZE, NUM_CLASSES);

    for (int epoch = 0; epoch < EPOCHS; epoch++) {
        float total_loss = 0.0f;

        for (int i = 0; i < num_train_samples; i++) {
            /* Receive activation and label */
            MPI_Recv(a2, HIDDEN2_SIZE, MPI_FLOAT, 1, TAG_FORWARD_A2, ...);
            MPI_Recv(&label, 1, MPI_INT, 0, TAG_FORWARD_LABEL, ...);

            /* Forward pass */
            dense_forward(&l3, a2, z3);
            softmax(z3, probs, NUM_CLASSES);

            /* Calculate loss */
            total_loss += cross_entropy_loss(probs, label);

            /* Backward pass */
            cross_entropy_backward(probs, label, grad_z3, NUM_CLASSES);
            dense_backward(&l3, grad_z3, grad_a2);

            /* Send gradient to Process 1 */
            MPI_Send(grad_a2, HIDDEN2_SIZE, MPI_FLOAT, 1, TAG_BACKWARD_GRAD2, ...);

            /* Update weights */
            sgd_update(&l3, LEARNING_RATE);
        }

        /* Send average loss to Process 0 */
        float avg_loss = total_loss / num_train_samples;
        MPI_Send(&avg_loss, 1, MPI_FLOAT, 0, TAG_LOSS, MPI_COMM_WORLD);
    }
}
\end{lstlisting}

% ============================================================
% DATA FLOW & CONTROL FLOW
% ============================================================
\section{Data Flow \& Control Flow}

\subsection{Data Flow Diagram}

\begin{figure}[h]
\centering
\begin{tcolorbox}[colback=white,colframe=black,width=\textwidth]
\begin{verbatim}
                    TRAINING: One Sample Iteration
    
    Process 0                Process 1                Process 2
    ─────────                ─────────                ─────────
    Load x, label
         │
    z1 = W1*x + b1
    a1 = ReLU(z1)
         │
         ├─── a1 (256 floats) ───►
         │                    Recv a1
         │                    z2 = W2*a1 + b2
         │                    a2 = ReLU(z2)
         │                         │
         │                         ├─── a2 (128 floats) ───►
         │                         │                    Recv a2
         ├─── label (1 int) ───────┼────────────────────►Recv label
         │                         │                    z3 = W3*a2 + b3
         │                         │                    probs = softmax(z3)
         │                         │                    loss = -log(probs[label])
         │                         │                    grad_z3 = probs - one_hot
         │                         │                    grad_a2 = W3^T * grad_z3
         │                         │                         │
         │                         ◄─── grad_a2 (128) ───────┤
         │                    grad_z2 = grad_a2 * ReLU'
         │                    grad_a1 = W2^T * grad_z2
         │                         │
         ◄─── grad_a1 (256) ───────┤
    grad_z1 = grad_a1 * ReLU'
    grad_W1, grad_b1 computed
    W1 -= lr * grad_W1
\end{verbatim}
\end{tcolorbox}
\caption{Complete data flow for one training sample}
\end{figure}

\subsection{Control Flow}

The control flow follows a strict \textbf{synchronous} pattern:

\begin{enumerate}
    \item All processes start and initialize their layers
    \item \texttt{MPI\_Bcast} synchronizes dataset sizes
    \item Training loop:
    \begin{itemize}
        \item Process 0 \textit{leads} by loading data and initiating forward pass
        \item Process 1 \textit{waits} for P0, then forwards to P2
        \item Process 2 \textit{waits} for P1, computes loss, \textit{initiates} backward pass
        \item Gradients flow back: P2 $\rightarrow$ P1 $\rightarrow$ P0
        \item All processes update weights independently
    \end{itemize}
    \item After each epoch, P2 sends loss to P0
    \item Evaluation follows similar forward-only pattern
    \item \texttt{MPI\_Finalize} ensures clean shutdown
\end{enumerate}

\subsection{Synchronization Points}

The program uses \textbf{blocking} MPI operations, creating implicit synchronization:

\begin{table}[h]
\centering
\begin{tabular}{@{}lll@{}}
\toprule
\textbf{Operation} & \textbf{Sender Blocks Until} & \textbf{Receiver Blocks Until} \\
\midrule
\texttt{MPI\_Send} & Buffer can be reused & N/A \\
\texttt{MPI\_Recv} & N/A & Message received \\
\texttt{MPI\_Bcast} & All processes call it & All processes call it \\
\bottomrule
\end{tabular}
\caption{MPI synchronization behavior}
\end{table}

% ============================================================
% ERROR HANDLING & EDGE CASES
% ============================================================
\section{Error Handling \& Edge Cases}

\subsection{Current Error Handling}

\begin{table}[h]
\centering
\begin{tabular}{@{}lll@{}}
\toprule
\textbf{Error Condition} & \textbf{Handling} & \textbf{Location} \\
\midrule
Wrong process count & Print error, exit & \texttt{main()} \\
File open failure & Print error, return NULL & \texttt{read\_idx\_*} \\
Dataset load failure & \texttt{MPI\_Abort} & \texttt{main()} \\
Log of zero & Clamp to \texttt{EPSILON} & \texttt{cross\_entropy\_loss} \\
\bottomrule
\end{tabular}
\caption{Error handling mechanisms}
\end{table}

\subsection{Edge Cases}

\begin{enumerate}
    \item \textbf{Empty dataset}: Would cause division by zero in loss averaging
    \item \textbf{Memory allocation failure}: \texttt{malloc}/\texttt{calloc} return values not checked
    \item \textbf{MPI failures}: MPI function return codes not checked
    \item \textbf{Numerical overflow}: Softmax uses max-subtraction to prevent overflow
    \item \textbf{All-negative ReLU outputs}: Could cause vanishing gradients (not handled)
\end{enumerate}

\subsection{Recommendations}

\begin{lstlisting}[language=C, caption={Recommended error checking}]
// Check malloc
float *buf = (float *)malloc(size * sizeof(float));
if (!buf) {
    fprintf(stderr, "Memory allocation failed\n");
    MPI_Abort(MPI_COMM_WORLD, 1);
}

// Check MPI operations
int ret = MPI_Send(...);
if (ret != MPI_SUCCESS) {
    char error_string[MPI_MAX_ERROR_STRING];
    int length;
    MPI_Error_string(ret, error_string, &length);
    fprintf(stderr, "MPI Error: %s\n", error_string);
    MPI_Abort(MPI_COMM_WORLD, 1);
}
\end{lstlisting}

% ============================================================
% COMPLEXITY & PERFORMANCE
% ============================================================
\section{Complexity \& Performance Notes}

\subsection{Time Complexity}

For one training sample with layer sizes $n_0 = 784$, $n_1 = 256$, $n_2 = 128$, $n_3 = 10$:

\begin{table}[h]
\centering
\begin{tabular}{@{}llc@{}}
\toprule
\textbf{Operation} & \textbf{Complexity} & \textbf{Dominant Term} \\
\midrule
Dense Forward (Layer 1) & $O(n_0 \times n_1)$ & $O(200,704)$ \\
Dense Forward (Layer 2) & $O(n_1 \times n_2)$ & $O(32,768)$ \\
Dense Forward (Layer 3) & $O(n_2 \times n_3)$ & $O(1,280)$ \\
Dense Backward (each) & Same as forward & -- \\
ReLU Forward/Backward & $O(n)$ & Negligible \\
Softmax & $O(n_3)$ & Negligible \\
\midrule
\textbf{Total per sample} & $O(n_0 n_1 + n_1 n_2 + n_2 n_3)$ & $\approx O(235,000)$ \\
\bottomrule
\end{tabular}
\caption{Time complexity analysis}
\end{table}

\subsection{Memory Complexity}

\begin{table}[h]
\centering
\begin{tabular}{@{}lll@{}}
\toprule
\textbf{Component} & \textbf{Size (floats)} & \textbf{Process} \\
\midrule
Layer 1 weights + gradients & $2 \times 784 \times 256 = 401,408$ & P0 \\
Layer 1 bias + gradients & $2 \times 256 = 512$ & P0 \\
Layer 2 weights + gradients & $2 \times 256 \times 128 = 65,536$ & P1 \\
Layer 3 weights + gradients & $2 \times 128 \times 10 = 2,560$ & P2 \\
Activation buffers & $\approx 2,000$ per process & All \\
Training images (P0 only) & $60,000 \times 784 = 47,040,000$ & P0 \\
\bottomrule
\end{tabular}
\caption{Memory usage breakdown}
\end{table}

\subsection{Communication Overhead}

Per training sample:
\begin{itemize}
    \item Forward: $256 + 128 + 1 = 385$ values $\approx 1.5$ KB
    \item Backward: $128 + 256 = 384$ values $\approx 1.5$ KB
    \item \textbf{Total}: $\approx 3$ KB per sample
\end{itemize}

For 60,000 samples $\times$ 10 epochs:
\begin{itemize}
    \item Total data transferred: $\approx 1.8$ GB
\end{itemize}

\subsection{Performance Bottlenecks}

\begin{enumerate}
    \item \textbf{Communication latency}: Blocking sends/receives create idle time
    \item \textbf{Pipeline bubble}: Processes wait for each other (no overlap)
    \item \textbf{Unbalanced workload}: Layer 1 has most computation
    \item \textbf{Sequential samples}: No batch-level parallelism
\end{enumerate}

\subsection{Potential Optimizations}

\begin{itemize}
    \item Use \texttt{MPI\_Isend}/\texttt{MPI\_Irecv} for non-blocking communication
    \item Implement \textbf{micro-batching} to overlap computation and communication
    \item Use \texttt{MPI\_Sendrecv} to combine operations
    \item Add SIMD vectorization for matrix operations
    \item Use BLAS libraries for dense operations
\end{itemize}

% ============================================================
% HOW TO RUN
% ============================================================
\section{How to Run / Example Usage}

\subsection{Prerequisites}

\begin{enumerate}
    \item Install MPI implementation:
    \begin{lstlisting}[language=bash]
# Ubuntu/Debian
sudo apt-get install openmpi-bin libopenmpi-dev

# macOS with Homebrew
brew install open-mpi

# CentOS/RHEL
sudo yum install openmpi openmpi-devel
    \end{lstlisting}
    
    \item Download MNIST dataset and place in \texttt{data/} directory
\end{enumerate}

\subsection{Compilation}

\begin{lstlisting}[language=bash, caption={Compilation commands}]
# Using Makefile
make

# Manual compilation
mpicc -Wall -O3 -std=c99 -o fcn_mpi fcn_mpi.c -lm
\end{lstlisting}

\subsection{Execution}

\begin{lstlisting}[language=bash, caption={Running the program}]
# Standard execution (requires 3+ CPU cores)
mpirun -np 3 ./fcn_mpi

# If you have fewer than 3 cores
mpirun -np 3 --oversubscribe ./fcn_mpi

# On a cluster with hostfile
mpirun -np 3 --hostfile hosts.txt ./fcn_mpi
\end{lstlisting}

\subsection{Expected Output}

\begin{lstlisting}[caption={Sample output}]
==============================================
Pipeline Parallelism Neural Network with MPI
==============================================
Architecture: 784 -> 256 -> 128 -> 10
Process 0: Input + Hidden Layer 1
Process 1: Hidden Layer 2
Process 2: Output Layer + Loss
==============================================

Loaded 60000 training samples and 10000 test samples

[Process 0] Starting training with 60000 samples
[Process 1] Ready for training
[Process 2] Ready for training
Epoch 1 | Loss: 0.5331 | Time: 17.35 s
Epoch 2 | Loss: 0.3920 | Time: 17.65 s
...
Epoch 10 | Loss: 0.2541 | Time: 18.51 s
[Process 0] Starting evaluation with 10000 samples
[Process 1] Ready for evaluation
[Process 2] Ready for evaluation
Test Accuracy: 84.92%

Total execution time: 189.37 seconds
\end{lstlisting}

% ============================================================
% TESTING SUGGESTIONS
% ============================================================
\section{Testing Suggestions}

\subsection{Unit Tests}

\begin{enumerate}
    \item \textbf{Tensor operations}:
    \begin{itemize}
        \item Verify \texttt{tensor\_create} allocates correct size
        \item Test \texttt{tensor\_zero} sets all elements to 0
        \item Check \texttt{tensor\_free} handles NULL gracefully
    \end{itemize}
    
    \item \textbf{Dense layer}:
    \begin{itemize}
        \item Forward pass with known weights (hand-calculated expected output)
        \item Backward pass gradient verification using finite differences
    \end{itemize}
    
    \item \textbf{Activations}:
    \begin{itemize}
        \item ReLU: Test with positive, negative, and zero values
        \item Softmax: Verify output sums to 1.0
        \item Softmax: Test numerical stability with large inputs
    \end{itemize}
    
    \item \textbf{Loss functions}:
    \begin{itemize}
        \item Cross-entropy with perfect prediction (should be $\approx 0$)
        \item Cross-entropy with zero probability (should use EPSILON)
    \end{itemize}
\end{enumerate}

\subsection{Integration Tests}

\begin{enumerate}
    \item \textbf{Single sample training}: Verify loss decreases
    \item \textbf{Overfitting test}: Train on 10 samples, expect $\sim$100\% accuracy
    \item \textbf{Determinism}: Same seed should produce same results
\end{enumerate}

\subsection{MPI-Specific Tests}

\begin{enumerate}
    \item \textbf{Message ordering}: Verify messages arrive in correct order
    \item \textbf{Deadlock detection}: Run with timeout to catch hangs
    \item \textbf{Process failure}: Test behavior if one process crashes
\end{enumerate}

\subsection{Performance Tests}

\begin{enumerate}
    \item \textbf{Scalability}: Compare with sequential version
    \item \textbf{Communication overhead}: Profile time spent in MPI calls
    \item \textbf{Memory usage}: Monitor with \texttt{valgrind} or \texttt{massif}
\end{enumerate}

% ============================================================
% IMPROVEMENTS
% ============================================================
\section{Improvements / Refactoring Opportunities}

\subsection{High Priority}

\begin{enumerate}
    \item \textbf{Add error checking for memory allocation}
    \begin{itemize}
        \item All \texttt{malloc}/\texttt{calloc} calls should check for NULL
        \item Use \texttt{MPI\_Abort} on failure for clean shutdown
    \end{itemize}
    
    \item \textbf{Check MPI function return values}
    \begin{itemize}
        \item All MPI calls can fail; check return codes
        \item Implement proper error handling and reporting
    \end{itemize}
    
    \item \textbf{Add batch processing}
    \begin{itemize}
        \item Current implementation processes one sample at a time
        \item Mini-batches would improve GPU utilization and convergence
    \end{itemize}
\end{enumerate}

\subsection{Medium Priority}

\begin{enumerate}
    \item \textbf{Use non-blocking MPI operations}
    \begin{lstlisting}[language=C]
MPI_Request req;
MPI_Isend(a1, size, MPI_FLOAT, 1, tag, comm, &req);
// Do computation while message is in flight
MPI_Wait(&req, &status);
    \end{lstlisting}
    
    \item \textbf{Implement micro-batching for pipeline efficiency}
    \begin{itemize}
        \item Process multiple samples simultaneously
        \item Reduces pipeline ``bubbles''
    \end{itemize}
    
    \item \textbf{Add model checkpointing}
    \begin{itemize}
        \item Save weights periodically
        \item Allow resuming training
    \end{itemize}
    
    \item \textbf{Configurable hyperparameters}
    \begin{itemize}
        \item Command-line arguments for learning rate, epochs, etc.
        \item Configuration file support
    \end{itemize}
\end{enumerate}

\subsection{Low Priority}

\begin{enumerate}
    \item \textbf{Use BLAS for matrix operations}
    \begin{itemize}
        \item Replace manual loops with optimized BLAS calls
        \item Significant performance improvement expected
    \end{itemize}
    
    \item \textbf{Add more layer types}
    \begin{itemize}
        \item Dropout for regularization
        \item Batch normalization
    \end{itemize}
    
    \item \textbf{Support variable number of processes}
    \begin{itemize}
        \item Dynamically assign layers to processes
        \item Better load balancing
    \end{itemize}
    
    \item \textbf{Add logging and profiling}
    \begin{itemize}
        \item Detailed timing information
        \item Communication statistics
    \end{itemize}
\end{enumerate}

\subsection{Code Quality}

\begin{enumerate}
    \item \textbf{Separate into multiple files}
    \begin{itemize}
        \item \texttt{tensor.c/h}, \texttt{layer.c/h}, \texttt{mnist.c/h}
        \item Improves maintainability and reusability
    \end{itemize}
    
    \item \textbf{Add comprehensive documentation}
    \begin{itemize}
        \item Doxygen-style comments for all functions
        \item README with architecture diagrams
    \end{itemize}
    
    \item \textbf{Implement proper cleanup on error}
    \begin{itemize}
        \item Use goto-based cleanup pattern
        \item Ensure all resources are freed on any exit path
    \end{itemize}
\end{enumerate}

% ============================================================
% OPEN QUESTIONS
% ============================================================
\section{Open Questions}

The following aspects are unclear or could not be determined from the code:

\begin{enumerate}
    \item \textbf{Weight initialization strategy}: Why use uniform $[-0.05, 0.05]$ instead of Xavier/He initialization?
    
    \item \textbf{Learning rate schedule}: Is a fixed learning rate optimal, or would decay improve results?
    
    \item \textbf{Batch size}: Why process samples one at a time instead of mini-batches?
    
    \item \textbf{Target accuracy}: What is the expected accuracy for this architecture on MNIST?
    
    \item \textbf{Scalability goals}: Is this intended for larger networks or just as a demonstration?
    
    \item \textbf{Dataset compatibility}: Does this work with Fashion-MNIST without modification?
    
    \item \textbf{Cluster deployment}: Are there specific MPI configurations recommended for multi-node execution?
\end{enumerate}

% ============================================================
% APPENDIX
% ============================================================
\appendix

\section{Complete Source Code}

The complete source code is available in \texttt{fcn\_mpi.c}. Key sections include:

\begin{itemize}
    \item Lines 1--38: Configuration and includes
    \item Lines 40--70: Tensor structure and operations
    \item Lines 72--145: Dense layer implementation
    \item Lines 147--200: Activation and loss functions
    \item Lines 202--270: MNIST loading
    \item Lines 272--350: Process 0 implementation
    \item Lines 352--410: Process 1 implementation
    \item Lines 412--480: Process 2 implementation
    \item Lines 482--540: Main function
\end{itemize}

\section{MPI Quick Reference}

\begin{table}[h]
\centering
\begin{tabular}{@{}ll@{}}
\toprule
\textbf{Function} & \textbf{Purpose} \\
\midrule
\texttt{MPI\_Init} & Initialize MPI environment \\
\texttt{MPI\_Finalize} & Clean up MPI environment \\
\texttt{MPI\_Comm\_rank} & Get process rank (ID) \\
\texttt{MPI\_Comm\_size} & Get total number of processes \\
\texttt{MPI\_Send} & Blocking send \\
\texttt{MPI\_Recv} & Blocking receive \\
\texttt{MPI\_Bcast} & Broadcast from one to all \\
\texttt{MPI\_Abort} & Terminate all processes \\
\texttt{MPI\_Wtime} & Wall-clock time \\
\bottomrule
\end{tabular}
\caption{MPI functions used in this program}
\end{table}

% ============================================================
% END DOCUMENT
% ============================================================
\end{document}
